\documentclass[10pt,a4paper]{amsart}
\usepackage[margin=19mm]{geometry}
\usepackage{amsmath,amssymb,mathtools}
\usepackage[scr=rsfs,cal=euler]{mathalfa}
\usepackage{xfrac}
\usepackage[unicode,hidelinks,bookmarks=false]{hyperref}
\newcommand{\ie}{i.e.\ }
\newcommand{\cf}{cf.\ }
\newcommand{\resp}{respectively\ }
\newcommand{\Subsection}{\subsection}
\providecommand{\nobreakdash}{\nobreak\hbox{-}}
\setlength{\emergencystretch}{2em}
\multlinegap0pt
\usepackage{bm}
\usepackage{bbm}
\usepackage{dsfont}
\usepackage{xspace}
\usepackage{cancel}
\usepackage{mathtools}
\usepackage{amsthm}
\makeatletter
\@ifundefined{theorem}{\newtheorem{theorem}{Theorem}}{}
\@ifundefined{thm}{\newtheorem{thm}[theorem]{Theorem}}{}
\makeatother
\def\Prob{{\mathbb P}}
\def\calM{\mathcal{M}}
\title[Resolving Memorization in Empirical Diffusion Model for Mani]{Resolving Memorization in Empirical Diffusion Model for Manifold Data in High-Dimensional Spaces}
\author{Yang Lyu, Tan Minh Nguyen, Yuchun Qian, Xin T. Tong}
\date{Source: arXiv:2505.02508v3 (CC BY 4.0)}
\begin{document}
\maketitle

\noindent\emph{Source and licence.} Resolving Memorization in Empirical Diffusion Model for Manifold Data in High-Dimensional Spaces. Version arXiv:2505.02508v3 (CC BY 4.0), distributed under the Creative Commons Attribution 4.0 International licence, as recorded in the arXiv licence metadata for this exact version and shown on the abstract page https://arxiv.org/abs/2505.02508v3. Full rights notice: licenses/arxiv-2505.02508-CC-BY-4.0.txt.\par\smallskip

\noindent\textit{Editorial scope.} Counted unit: the Thm whose source label is \texttt{thm:NWproj} in the article (source0.tex lines 638--653, source label \texttt{thm:NWproj}), delivered together with the article's own complete original proof of it (source0.tex lines 655--788, one \texttt{proof} environment of 134 lines). No mathematics is written, repaired or re-derived: statement and proof are the article's own text line for line. Required context retained uncounted: none. Every definition, display and label of those lines is retained unchanged. Omitted from this package and not redistributed: 1--637, 654--654, 789--2098. Nothing inside a retained excerpt is omitted or altered: every retained body is delivered exactly as the article's lines are written, so no omission receipt of any kind accompanies this package. Citations: the retained excerpts carry no citation command at all, so no bibliography entry of the article is transcribed and no reference is redistributed. Reference adaptation: none was needed and none was made; every reference inside the delivered unit resolves inside it, so no unresolved reference or citation is produced. Printed slips kept literal: no printed word, symbol, hypothesis or number of the counted statement or of its proof was altered. Layout and class: the article's own class is \texttt{article} with packages including xinstyle, algorithm, algpseudocode, inputenc, fontenc, hyperref, url, booktabs, amsfonts, nicefrac, microtype, xcolor, amsmath, amssymb; the wrapper is the lane's pinned \texttt{amsart} editorial wrapper at 10pt on A4, which restates the theorem-like environments the retained text needs and adds the article's own macro definitions that the retained text needs (\texttt{\textbackslash Prob}, \texttt{\textbackslash calM}), with extra wrapper packages (bm, bbm, dsfont, xspace, cancel, mathtools). The delivered numbering is the unit's own. Assets: no retained span contains a figure environment, a graphics-inclusion command, a PDF-page inclusion, a table or a TikZ picture; no image file is copied into this package, the package declares no asset dependency and no image file is redistributed with it. Licence: version arXiv:2505.02508v3 of this article is distributed under the Creative Commons Attribution 4.0 International licence, as recorded in the arXiv licence metadata for this exact version and shown on the abstract page https://arxiv.org/abs/2505.02508v3. A full rights notice is delivered at licenses/arxiv-2505.02508-CC-BY-4.0.txt. Characters: every retained excerpt is pure ASCII, so the delivered engine, XeLaTeX with its default Latin Modern text font, reports no missing character.
\begin{thm}
\label{thm:NWproj}
Given any fixed realization $X=[X_{(1)},\ldots,X_{(n)}]$ and $x\in \calM$ and $r>0$. Suppose 
\[
\inf_{y\in\calM}|I_{y,\sigma}|\geq  1, 
\]
and $n,\sigma^{-1}$ are sufficiently large.
Let $\xi\sim \mathcal{N}(0,I)$ with decomposition $\xi=\xi_{\mathcal{T}}+\xi_\bot$ on $T_x\calM\oplus N_x\calM$, then 
\[
\mathbb{P}\left( \left\| F_{X,\sigma}(x + \sigma \xi) - F_{X,\sigma}(\exp_x(\sigma \xi_{\mathcal{T}}))\right\| \lesssim (d(k+1)\log n\, \sigma)^2|X \right) \ge 1 - n^{-k}. 
\]
In particular, 
\[
\mathbb{P}\left( \left\| F_{X,\sigma}(x + \sigma \xi_\bot) - F_{X,\sigma}(x)\right\| \lesssim (d(k+1)\log n\, \sigma)^2|X \right) \ge 1 - n^{-k}. 
\]
\end{thm}

\begin{proof}
\textbf{Step 1: A Pythagorean approximation.}
We start by showing for $r=\sqrt{2d(k+1)\log n}\sigma,$ 
\begin{align}
\notag
\mathbb{P}\bigg(  \big|\|x + \sigma \xi &- X_{(i)}\|^2 - \|x + \sigma \xi_{\mathcal{T}} - X_{(i)}\|^2\\ &- \|\sigma \xi_\bot\|^2 \big| \lesssim \sqrt{k \log n} \, \sigma r^2,\ 
\forall i \in I_{x,r}|X\bigg) \ge 1 - n^{-k}.
\label{tmp:orthogonal}
\end{align}
We decompose
\begin{align}
\notag
\Delta_{1,i}:&=\|x + \sigma \xi - X_{(i)}\|^2 - \|x + \sigma \xi_{\mathcal{T}} - X_{(i)}\|^2 - \|\sigma \xi_\bot\|^2 \\
\label{tmp:decomp}
&= 2 \sigma \langle x + \sigma \xi_{\mathcal{T}} - X_{(i)}, \xi_\bot \rangle  = 2 \sigma \langle x - X_{(i)}, \xi_\bot \rangle.     
\end{align}
Since $i \in I_{X,r}$, $\exists v_i \in T_x \mathcal{M}$ with $\|v_i\|\leq 2r$ so that,  
\[
X_{(i)} = \exp_x(v_i) = x + v_i + O(\|v_i\|^2).
\]
This indicates, 
\[
\langle x - X_{(i)}, \xi_\bot \rangle = \langle x + v_i - \exp_x(v_i), \xi_\bot \rangle.
\]
Note that
\[
\|x + v_i - \exp_x(v_i)\| = \mathcal{O}(\|v_i\|^2) = \mathcal{O}(r^2).
\]
Also, $\langle x + v_i - \exp_x(v_i), \xi_\bot \rangle \sim \mathcal{N}(0, \|P_\bot(x + v_i - \exp_x(v_i))\|^2)$. Here, $P_\bot: \mathbb{R}^D \to N_x \mathcal{M}$ is the orthogonal projection onto $N_x \mathcal{M}$.\\

Let $t_i := \|P_\bot(x + v_i - \exp_x(v_i))\| = \mathcal{O}(r^2)$, $q_i := \langle x + v_i - \exp_x(v_i), \xi_\bot \rangle$, standard Gaussian tail inequality says
\[
\mathbb{P}(|q_i| > z|X) \le 2 \exp\left( -\frac{z^2}{2 t_i^2} \right).
\]
Apply union bound for the event $\mathcal{A}=\{\max_{i \in I_{X,r}} |q_i| > z\}$, then for some constant $C$
\[
\mathbb{P}\left( \mathcal{A} \,|\,X\right) \leq 2 |I_{X,r}|\exp\left( -\frac{z^2}{C r^2} \right).
\]
To have 
\[
2 n\exp\left( -\frac{z^2}{2 C r^4} \right) \leq n^{-k},
\]
It suffices to have $z \gtrsim r^2 \sqrt{(k+1) \log n}.$ Using \eqref{tmp:decomp}, we have \eqref{tmp:orthogonal}.


\noindent\textbf{Step 2: Show $\|\exp_x(\sigma \xi_{\mathcal{T}})-X_{(i)}\|^2-\|x+\sigma \xi_{\mathcal{T}}-X_{(i)}\|^2=\widetilde{O}_p(\sigma^3)$ for $i\in I_{X,r}$.} 

Denote the event $\mathcal{B}=\{\|\xi_{\mathcal{T}}\|\geq \sqrt{d(k+1)\log n }\}$. By Gaussian tail probability, we know $\Prob(\mathcal{B}|X)\leq n^{-k}$. Note under $\mathcal{B}^c$, for $i\in I_{x,r}$
\begin{align*}
|\Delta_{2,i}|&:=|\|\exp_x(\sigma \xi_{\mathcal{T}})-X_{(i)}\|^2-\|x+\sigma \xi_{\mathcal{T}}-X_{(i)}\|^2|\\
&\leq (\|\exp_x(\sigma \xi_{\mathcal{T}})-X_{(i)}\|+\|x+\sigma \xi_{\mathcal{T}}-X_{(i)}\|)\|\exp_x(\sigma \xi_{\mathcal{T}})-x-\sigma\xi_{\mathcal{T}}\|\\
&\leq 4Cr\|\sigma \xi_{\mathcal{T}}\|^2\leq 4C \sigma^3 (d(k+1)\log n )^{3/2}. 
\end{align*}




\noindent\textbf{Step 3: Showing $F_{X,\sigma}(x+\sigma\xi)=F_{X,\sigma}(\exp_x(\sigma \xi_{\mathcal{T}}))+\widetilde{O}_p(\sigma^2)$.}
Denote $x'=\exp_x(\sigma \xi_{\mathcal{T}})$, we note that if we let $z_i=\Delta_{1,i}+\Delta_{2,i}$
\begin{align*}
k_\sigma(x+\sigma \xi, X_{(i)})&=\exp(-\frac1{2\sigma^2}\|x+\sigma\xi-X_{(i)}\|^2)\\
&=\exp(-\frac1{2\sigma^2}\|x'-X_{(i)}\|^2-\frac1{2}\|\xi_\bot\|^2)\exp(-\frac{z_i}{2\sigma^2})\\
&=\exp(-\frac1{2}\|\xi_\bot\|^2)\exp(-\frac{z_i}{2\sigma^2})k_\sigma(x', X_{(i)}).
\end{align*}
Therefore we have
\begin{align*}
F_{X, \sigma}(x')= \frac{\sum\limits_{i=1}^n k_\sigma(x', X_{(i)}) X_{(i)}}{\sum\limits_{i=1}^n k_\sigma(x', X_{(i)})}
=\frac{\sum_{i\in [n]} w_i X_{(i)}}{\sum_{i\in [n]} w_i},
\end{align*}
and 
\begin{align*}
F_{X, \sigma}(x+\sigma \xi)=
\frac{\sum\limits_{i=1}^n k_\sigma(x', X_{(i)})\exp(-\frac{z_i}{2\sigma^2}) X_{(i)}}{\sum\limits_{i=1}^n k_\sigma(x', X_{(i)})\exp(-\frac{z_i}{2\sigma^2})}
=\frac{\sum_{i\in [n]} w_iu_i X_{(i)}}{\sum_{i\in [n]} w_iu_i},
\end{align*}
where $w_i:=k_\sigma(x', X_{(i)})$ and $u_i:=\exp(-\frac{z_i}{2\sigma^2})$. 
So 
\begin{align*}
&\left\|\frac{\sum_{i\in [n]} w_iu_i X_{(i)}}{\sum_{i\in [n]} w_iu_i}-
\frac{\sum_{i\in [n]} w_i X_{(i)}}{\sum_{i\in [n]} w_i}\right\|\\
&=\left\|\frac{\sum_{i\in [n]} w_iu_i (X_{(i)}-x')}{\sum_{i\in [n]} w_iu_i}-
\frac{\sum_{i\in [n]} w_i (X_{(i)}-x')}{\sum_{i\in [n]} w_i}\right\|\\
&\leq\left\|\frac{\sum_{i\in [n]} w_i (X_{(i)}-x')}{\sum_{i\in [n]} w_iu_i}-\frac{\sum_{i\in [n]} w_i (X_{(i)}-x')}{\sum_{i\in [n]} w_i}\right\|+\left\|\frac{\sum_{i\in [n]} w_i(u_i-1) (X_{(i)}-x')}{\sum_{i\in [n]} w_iu_i}\right\|,
\end{align*}
% Denote the event $\mathcal{B}=\{\|\xi_{\mathcal{T}}\|>\sqrt{2d(k+1)\log n}\}$. 
% Because $\xi_{\mathcal{T}}\sim\mathcal{N}(0,d)$, we know  $\Prob(\mathcal{B}|X)\leq  n^{-k}$. 
% Under $\mathcal{B}$ 
 Note that if $i\notin I_{x',r}$ with $r=\sqrt{2d(k+1)\log n}\sigma$  
\[
w_i=\exp(-\frac{1}{2\sigma^2} \|x'-X_{(i)}\|^2)
\leq \exp(-(k+1) \log n)\leq n^{-k-1}. 
\]
So 
\[
\sum_{i\notin I_{x',r}} w_i\leq n^{-k},\quad
\sum_{i\notin I_{x',r}} w_iu_i\leq n^{-k},\quad
\left\|\sum_{i\notin I_{x',r}} w_i (X_{(i)}-x')\right\|\leq n^{-k}\text{diam}(\calM). 
\]

Note that for all $i\in I_{x,r}$, under $\mathcal{A}^c\cap \mathcal{B}^c$, 
\[
\max\{|u_i-1|,\frac{|u_i-1|}{u_i}\}\leq \frac{|z_i|}{2\sigma^2}\leq C\sigma (d(k+1)\log n)^{3/2}<1/2,
\]
where we use the exponential ratio bound: 
for all \( x \in \mathbb{R} \) with small enough norm, we have:
\[
\left| e^{-x} - 1 \right| \le |x|.
\]
Meanwhile by our assumption of $X$,
\[
\sum_{i}w_i\geq  \sum_{i\in I_{x',\sigma}}w_i\geq 
\sum_{i\in I_{x',\sigma}}\exp(-1/2)
\geq e^{-1/2}.
\]
Meanwhile under $\mathcal{A}^c\cap \mathcal{B}^c$, $I_{x',\sigma}\subset I_{x,r}$, $\sum_i u_i w_i\geq 
\sum_{i\in I_{x',\sigma}}\exp(-1/2) u_i\geq \frac12 e^{-1/2}$. 

We bound each term using 
\begin{align*}
&\left\|\frac{\sum_{i\in [n]} w_i (X_{(i)}-x')}{\sum_{i\in [n]} w_iu_i}-\frac{\sum_{i\in [n]} w_i(X_{(i)}-x')}{\sum_{i\in [n]} w_i}\right\|\\
&=\frac{\|\sum_{i\in [n]} w_i (X_{(i)}-x')\|(\sum_{i\in [n]}w_i u_i|u_i^{-1}-1|)}{(\sum_{i\in [n]}w_i u_i)(\sum_{i\in [n]}w_i)}\\
&\leq \frac{\|\sum_{i\in [n]} w_i (X_{(i)}-x')\|(\sum_{i\in [n]}w_iu_i)\sigma (d(k+1)\log n)^{3/2}}{(\sum_{i\in [n]}w_i u_i)(\sum_{i\in [n]}w_i)}\\
&\leq \frac{(\|\sum_{i\in I_{x',r}} w_i (X_{(i)}-x')\|+\text{diam}(\calM)n^{-k})}{\sum_{i\in [n]}w_i}\sigma (d(k+1)\log n)^{3/2}\\
&\leq \left(r+C_d\text{diam}(\calM)n^{-k}\right)\sigma (d(k+1)\log n)^{3/2}\\
&\leq \sigma^2  (d(k+1)\log n)^2. 
\end{align*}
Similarly
\begin{align*}
&\left\|\frac{\sum_{i\in [n]} w_i(u_i-1) (X_{(i)}-x)}{\sum_{i\in [n]} w_iu_i}\right\|\\
&\leq \left\|\frac{\sum_{i\in I_{x,r}} w_i(u_i-1) (X_{(i)}-x)}{\sum_{i\in [n]} w_iu_i}\right\|+\left\|\frac{\sum_{i\in I^c_{x,r}} w_i(u_i-1) (X_{(i)}-x)}{\sum_{i\in [n]} w_iu_i}\right\|\\
&\leq r\left|\frac{\sum_{i\in I_{x,r}} w_i(u_i-1) }{\sum_{i\in [n]} w_iu_i}\right|+\frac{n^{-k}\text{diam}(\calM)}{\sum_{i\in [n]} w_iu_i}\leq 2e\sigma^2  ((k+1)\log n)^2.
\end{align*}
Finally, note that $\xi_{\mathcal{T}}$ and $\xi_\bot$ are independent. So all the discussion above holds if we add in the condition that $\xi_{\mathcal{T}}=0$. 
\end{proof}

\end{document}
